% In this file you should put macros to be used only by
% the printed version. Of course they should have a corresponding
% version in macros/web.tex.
% Typically the printed version could have more fancy decorations.
% This should be a very short file.
%
% This file starts with dummy macros that ensure the pdf
% compiler will ignore macros provided by plasTeX that make
% sense only for the web version, such as dependency graph
% macros.


% Dummy macros that make sense only for web version.
\newcommand{\lean}[1]{}
\newcommand{\discussion}[1]{}
\newcommand{\leanok}{}
\newcommand{\mathlibok}{}
\newcommand{\notready}{}
% Make sure that arguments of \uses and \proves are real labels, by using invisible refs:
% latex prints a warning if the label is not defined, but nothing is shown in the pdf file.
% It uses LaTeX3 programming, this is why we use the expl3 package.
\ExplSyntaxOn
\NewDocumentCommand{\uses}{m}
 {\clist_map_inline:nn{#1}{\vphantom{\ref{##1}}}%
  \ignorespaces}
\NewDocumentCommand{\proves}{m}
 {\clist_map_inline:nn{#1}{\vphantom{\ref{##1}}}%
  \ignorespaces}
\ExplSyntaxOff

% Cross-reference into the Bestiary chapter: in print this is a plain
% hyperref to the label. The web version (macros/web.tex) emits a direct
% chap-bestiary.html anchor link instead, since plasTeX's \hyperref is broken.
\newcommand{\bestiaryref}[2]{\hyperref[#1]{#2}}

% In the printed version we just render the declaration name in
% typewriter --- the URL would clutter the page. Web version provides a
% live link via macros/web.tex.
% A Lean name may break after any `.` or `_` (no hyphen: a hyphen would read as part
% of the name), so a long \leandocref at a line end stays inside the margin.
\ExplSyntaxOn
\tl_new:N \l__leandoc_tl
\NewDocumentCommand{\leandocref}{m}{
  \tl_set:Nn \l__leandoc_tl {#1}
  \tl_replace_all:Nnn \l__leandoc_tl {.} {.\allowbreak}
  \tl_replace_all:Nnn \l__leandoc_tl {\_} {\_\allowbreak}
  \texttt{\l__leandoc_tl}
}
\ExplSyntaxOff

% \prosesection: a chapter-body prose sub-heading. In the printed version
% it is a normal unnumbered section, exactly as before. The web version
% (macros/web.tex) demotes it below the web TOC depth so the ch1-9 web
% sidebars match the PDF table of contents (numbered sections only).
\newcommand{\prosesection}[1]{\section*{#1}}
% \proseparagraph: see macros/web.tex; in print the ordinary run-in heading.
\newcommand{\proseparagraph}[1]{\paragraph{#1}}

% \cellleft: ragged-right text in a p{...} tabular cell. \raggedright alone
% turns the row-ending \\ into a line break, so the row end is restored.
% The web version (macros/web.tex) is empty: plasTeX sets cells as plain
% left-aligned HTML, and its \raggedright wraps the cell in a block that
% breaks the inline text apart.
\newcommand{\cellleft}{\raggedright\let\\\tabularnewline}

% \setbookmarksdepth: the PDF bookmark depth (hyperref). The web version
% (macros/web.tex) is empty: plasTeX has no \hypersetup and would print
% its argument as page text.
\newcommand{\setbookmarksdepth}[1]{\hypersetup{bookmarksdepth=#1}}
